\section{Tema 5}
\section*{1. Delimitación y clasificación}

\begin{enumerate}[label=\textbf{\arabic*.}]

    \item \textbf{Los servicios no destinados a la venta se definen por suministrarse a través del mercado a precios que cubren íntegramente sus costes de producción.} \\
    \textbf{Falsa.} \textit{Justificación: Los servicios no destinados a la venta son aquellos suministrados por el sector público, ya sea de forma completamente gratuita o a un precio que está desvinculado de su coste real de producción.}

    \item \textbf{Los denominados ``servicios estancados'' son aquellos que permiten una creciente capitalización e incorporación de tecnología para mejorar su productividad.} \\
    \textbf{Falsa.} \textit{Justificación: Los ``servicios estancados'' son precisamente los que muestran grandes dificultades para reducir la mano de obra por unidad de producto. Son los denominados ``servicios progresivos'' los que permiten una creciente capitalización y la incorporación de mejoras tecnológicas.}

    \item \textbf{Una de las principales dificultades para medir la productividad en servicios como la Administración Pública es que la producción se estima a menudo mediante los recursos o ``inputs'' utilizados.} \\
    \textbf{Verdadera.} \textit{Justificación: Dado que el output es difícil de cuantificar objetivamente, a menudo se mide por los costes o recursos empleados, complicando el análisis real de la productividad.}

    \item \textbf{Las actividades de contabilidad o publicidad realizadas internamente por una empresa industrial se contabilizan estadísticamente como parte del valor añadido del sector servicios.} \\
    \textbf{Falsa.} \textit{Justificación: Los servicios que producen internamente las empresas de otros sectores (como la industria o la agricultura) no se contabilizan dentro del sector terciario. Únicamente se imputan al sector servicios cuando estas actividades se externalizan a empresas especializadas.}

    \item \textbf{La necesaria concurrencia de la producción y el consumo en el espacio y el tiempo es un rasgo característico de gran parte de las actividades terciarias.} \\
    \textbf{Verdadera.} \textit{Justificación: La intangibilidad y la simultaneidad temporal y espacial entre la prestación y el consumo son características intrínsecas de los servicios.}

\end{enumerate}

\section*{2. Evolución del sector}

\begin{enumerate}[label=\textbf{\arabic*.}]
    \setcounter{enumi}{5}

    \item \textbf{El proceso de terciarización de la economía española se produjo de forma temprana y simultánea al de los principales países de la Unión Europea.} \\
    \textbf{Falsa.} \textit{Justificación: El proceso de terciarización en España ha sido un fenómeno reciente que comenzó a intensificarse fundamentalmente a partir de los años ochenta. En las principales economías de la Unión Europea, este proceso se produjo de forma anterior.}

    \item \textbf{El concepto de ``sector refugio'' en los servicios se apoya en factores de oferta tradicionales como el predominio de las PYMEs y los bajos costes de acceso al mercado.} \\
    \textbf{Verdadera.} \textit{Justificación: Efectivamente, ciertas ramas del sector servicios actúan como ``sector refugio'' porque presentan muy pocas barreras de entrada, predominan empresas de reducida dimensión (como en la hostelería y pequeña comercialización) y requieren una baja inversión inicial.}

    \item \textbf{La expansión del peso de los servicios en la producción nacional a precios corrientes se explica fundamentalmente porque sus precios han crecido a un ritmo inferior al de los bienes industriales.} \\
    \textbf{Falsa.} \textit{Justificación: Es justo al contrario. La participación de los servicios ha aumentado en términos nominales (precios corrientes) precisamente porque sus precios se han encarecido de forma relativa a un ritmo muy superior al de los bienes industriales.}

    \item \textbf{El sector servicios tiende a aumentar su participación relativa en la producción real de la economía durante las etapas de crisis económica.} \\
    \textbf{Verdadera.} \textit{Justificación: Durante las crisis, los sectores como la industria o la construcción suelen sufrir desplomes mucho más pronunciados en su demanda, por lo que el sector servicios retiene o incluso aumenta mecánicamente su participación relativa en el conjunto de la producción real de la economía.}

    \item \textbf{El desplazamiento masivo del empleo hacia el sector servicios en España es una consecuencia directa del crecimiento superior de su productividad frente a la agricultura o la industria.} \\
    \textbf{Falsa.} \textit{Justificación: El incremento del empleo en el sector servicios obedece a un cambio estructural derivado de la sustitución de mano de obra por otros factores (maquinaria, tecnología) en la agricultura y la industria. De hecho, la productividad en los servicios ha registrado un avance muy lento, no superior.}

    \item \textbf{Históricamente, la gran mayoría de las actividades terciarias en España han estado plenamente expuestas a la competencia internacional, reflejando su peso en el empleo.} \\
    \textbf{Falsa.} \textit{Justificación: Históricamente, y hasta el año 1990, la inmensa mayoría de los servicios en España se desarrollaron completamente al margen de las relaciones exteriores. El turismo era la única rama plenamente expuesta a la competencia internacional. Aún hoy, numerosas actividades presentan nula o escasa competencia exterior.}

\end{enumerate}

\section*{3. Especialización productiva y comercial}

\begin{enumerate}[label=\textbf{\arabic*.}]
    \setcounter{enumi}{11}

    \item \textbf{El fuerte avance de los servicios no destinados a la venta en España se justifica exclusivamente por la asunción de tareas educativas y sanitarias por parte del Estado.} \\
    \textbf{Falsa.} \textit{Justificación: Aunque la asunción de tareas educativas, sanitarias y asistenciales (el Estado de Bienestar) es clave, no es la única causa. Este avance también se justifica por la profunda descentralización de las Administraciones Públicas y por el encarecimiento relativo de la prestación de dichos servicios.}

    \item \textbf{Dentro de la estructura productiva española, la rama de hostelería presenta una relevancia significativamente mayor que en el promedio de la Unión Europea.} \\
    \textbf{Verdadera.} \textit{Justificación: La hostelería representa casi el triple de peso en la estructura productiva española en comparación con la media de sus homólogos europeos.}

    \item \textbf{Los servicios vinculados al bienestar, como la sanidad y la educación, mantienen en España un peso relativo en la producción superior a la media de la UE-27.} \\
    \textbf{Falsa.} \textit{Justificación: Los datos demuestran lo contrario. Las actividades sanitarias y de servicios sociales representan un 8,4\% en España frente al 10,2\% de la UE-27. La educación supone el 6,7\% frente al 6,6\%, situándose en niveles inferiores o iguales a la media europea.}

    \item \textbf{En la estructura del comercio exterior español, la actividad de ``turismo y viajes'' es la única que registra un superávit comercial en la balanza de servicios.} \\
    \textbf{Falsa.} \textit{Justificación: Si bien el turismo es un pilar, no es el único. Otras ramas de servicios aportan un fuerte superávit con altas tasas de cobertura, como los Servicios Financieros (308,8\%), Telecomunicaciones o Transportes.}

    \item \textbf{Aunque el turismo sigue siendo el pilar fundamental, las exportaciones españolas de servicios han mostrado una tendencia hacia la diversificación en las últimas décadas.} \\
    \textbf{Verdadera.} \textit{Justificación: Se aprecia un creciente empuje y exportación en sectores no turísticos, aportando diversidad a la balanza exterior.}

\end{enumerate}

\section*{4. Eficiencia productiva}

\begin{enumerate}[label=\textbf{\arabic*.}]
    \setcounter{enumi}{16}

    \item \textbf{Las ganancias de productividad han sido el motor principal del crecimiento de la producción total de servicios en España durante las últimas décadas.} \\
    \textbf{Falsa.} \textit{Justificación: El avance de la productividad en el sector terciario ha sido sumamente lento. Su aportación al crecimiento de la producción total fue de apenas un 29,8\% en España, inferior al promedio europeo, apoyándose el sector más en la acumulación de empleo que en mejoras de eficiencia.}

    \item \textbf{La incorporación de las nuevas tecnologías de la información ha permitido ganancias de productividad en la distribución comercial a través de la expansión de las ventas online.} \\
    \textbf{Verdadera.} \textit{Justificación: La digitalización y el e-commerce han supuesto un vector claro de eficiencia para el segmento de la distribución.}

    \item \textbf{El descenso de la productividad en la hostelería española se explica en parte por la gran dimensión de sus empresas, que genera ineficiencias organizativas.} \\
    \textbf{Falsa.} \textit{Justificación: Se explica justo por lo contrario: por la reducida dimensión media de sus empresas, lo cual les impide aprovechar las economías de escala y optimizar los costes operativos.}

    \item \textbf{España presenta una posición de productividad por hora trabajada superior a la media de la UE-27 en actividades financieras y de seguros.} \\
    \textbf{Verdadera.} \textit{Justificación: Es uno de los pocos subsectores donde España goza de una posición productiva ampliamente superior al promedio comunitario.}

    \item \textbf{El lento avance de la productividad en los servicios en España se debe a una especialización en actividades con elevada intensidad de capital físico.} \\
    \textbf{Falsa.} \textit{Justificación: El bajo crecimiento productivo es consecuencia de la especialización de España en actividades caracterizadas por una baja intensidad en capital físico y humano, junto con una elevada intensidad en mano de obra con baja cualificación.}

    \item \textbf{La escasa competencia en numerosas actividades terciarias ha actuado como un freno para la modernización tecnológica y organizativa del sector en España.} \\
    \textbf{Verdadera.} \textit{Justificación: Las rigideces y la falta de presión competitiva retrasan las inversiones en modernización y mejora organizativa.}

\end{enumerate}

\section*{5. Política sectorial}

\begin{enumerate}[label=\textbf{\arabic*.}]
    \setcounter{enumi}{22}

    \item \textbf{La liberalización de servicios en Europa desde los años 90 ha implicado una mayor intervención en la fijación de precios por parte del Estado.} \\
    \textbf{Falsa.} \textit{Justificación: El proceso de liberalización ha implicado exactamente lo contrario: la eliminación de controles estatales y la consecuente libertad de fijación de precios por las fuerzas del mercado.}

    \item \textbf{La política de liberalización de los servicios en España ha generado una redistribución de la renta desde los consumidores hacia los empresarios y trabajadores del sector.} \\
    \textbf{Falsa.} \textit{Justificación: El flujo distributivo ha sido a la inversa. Al aumentar la competencia y caer los precios de los servicios, la liberalización ha generado una redistribución de la renta desde los empresarios y trabajadores del sector, hacia los consumidores finales.}

    \item \textbf{En el transporte ferroviario español de pasajeros, la legislación actual impide la entrada de empresas privadas para competir con el operador público en alta velocidad.} \\
    \textbf{Falsa.} \textit{Justificación: La legislación actual ha abierto el mercado y permite plenamente la competencia. De hecho, junto a Renfe, operan actualmente tres compañías privadas con licencia y redes propias en alta velocidad (como Ouigo e Iryo).}

    \item \textbf{La liberalización del transporte aéreo en España ha permitido mejorar la accesibilidad y conectividad de los consumidores gracias a la entrada de aerolíneas de bajo coste.} \\
    \textbf{Verdadera.} \textit{Justificación: Ha democratizado los vuelos, multiplicando las rutas y abaratando los precios.}

    \item \textbf{En el sector de la distribución comercial, la existencia de normativas diferentes entre las regiones españolas supone un obstáculo para el desarrollo de estrategias globales de empresa.} \\
    \textbf{Verdadera.} \textit{Justificación: La heterogeneidad y fragmentación legislativa entre Comunidades Autónomas encarece y dificulta la expansión de las grandes distribuidoras a nivel nacional.}

    \item \textbf{A pesar del proceso de apertura a la competencia, el mercado de telefonía móvil en España muestra una tendencia reciente hacia la fragmentación, reduciendo la cuota de mercado de los grandes operadores.} \\
    \textbf{Falsa.} \textit{Justificación: Tras diversas fusiones y compras corporativas recientes, el mercado de telefonía móvil muestra un aumento de la concentración, donde los tres grandes operadores controlan más del 90\% de la cuota de mercado.}

    \item \textbf{En la comercialización de productos farmacéuticos, la normativa vigente prohíbe explícitamente que las farmacias desarrollen estrategias de integración vertical u horizontal.} \\
    \textbf{Verdadera.} \textit{Justificación: Es una de las barreras y protecciones estructurales que impiden a las farmacias formar redes, cadenas corporativas o fusionarse con las distribuidoras.}

    \item \textbf{La entrada de nuevos competidores en el mercado postal ha desplazado a la Sociedad Estatal Correos y Telégrafos, que ha dejado de ser el operador líder del sector.} \\
    \textbf{Falsa.} \textit{Justificación: Aunque la competencia se ha intensificado notablemente en el sector postal, empujando a mejoras operativas, el operador público (Correos) sigue manteniendo su posición como líder hegemónico del mercado.}

\end{enumerate}

